% !TEX program = pdflatex
% ------------------------------------------------------------
% MATH 431 Homework Template (plug-and-chug friendly)
%
% Goals:
% - Same clean header style as your MATH 521 template.
% - Easy “Problem -> Solution” workflow.
% - Lightweight (no tcolorbox); uses mdframed for optional boxes.
% - Good for probability homework: sample space, probability measure, events.
%
% Quick use:
%   1) Edit the Metadata block.
%   2) For each exercise, use \problem{<section>.<number>} and paste the prompt.
%   3) Write your solution in the solution environment.
% ------------------------------------------------------------

\documentclass[12pt]{article}

% =========================
% 0) Metadata (EDIT ME)
% =========================
\newcommand{\CourseCode}{MATH 431}
\newcommand{\CourseTitle}{Introduction to Probability}
\newcommand{\CourseSection}{LEC 002} % change to 001 or 003 if needed
\newcommand{\Semester}{Spring 2026}
\newcommand{\Instructor}{Sarah Tammen}
\newcommand{\HomeworkNumber}{8}
\newcommand{\YourName}{Alejandro De La Torre}
\newcommand{\DueDate}{March 20, 2026} 
% =========================
% 1) Core math tools
% =========================
\usepackage{amsmath, amssymb, amsthm}

% =========================
% 2) Lists and spacing
% =========================
\usepackage{enumitem}
\setlist[enumerate,1]{label=\textbf{(\alph*)}, leftmargin=2em, itemsep=0.25em}
\setlist[itemize]{leftmargin=1.5em, itemsep=0.25em}
\setlength{\parskip}{0.35em}
\setlength{\parindent}{0pt}

% =========================
% 3) Page geometry + header/footer
% =========================
\usepackage{geometry}
\geometry{margin=1in}

\usepackage{fancyhdr}
\pagestyle{fancy}
\fancyhf{}
\lhead{\CourseCode\ --- \CourseSection, \Semester}
\rhead{Homework \HomeworkNumber}
\cfoot{\thepage}
\renewcommand{\headrulewidth}{0.4pt}
\setlength{\headheight}{15pt}
\addtolength{\topmargin}{-2pt}

% =========================
% 4) Links (subtle)
% =========================
\usepackage[hidelinks]{hyperref}
\setcounter{tocdepth}{1}

% =========================
% 5) Figures (optional)
% =========================
\usepackage{graphicx}
\graphicspath{{images/}}

% =========================
% 6) Lightweight boxes (optional)
% =========================
\usepackage{mdframed}
\newmdenv[
  skipabove=0.6\baselineskip,
  skipbelow=0.6\baselineskip,
  linewidth=0.6pt,
  roundcorner=4pt,
  innerleftmargin=10pt,
  innerrightmargin=10pt,
  innertopmargin=8pt,
  innerbottommargin=8pt
]{boxnote}

% =========================
% 7) Probability / notation macros
% =========================
\newcommand{\R}{\mathbb{R}}
\newcommand{\N}{\mathbb{N}}
\newcommand{\Z}{\mathbb{Z}}

\newcommand{\OmegaSpace}{\Omega}
\newcommand{\FField}{\mathcal{F}}
% Probability measure in case it is relevant or want to distinguish it from P(A)
%\newcommand{\PMeasure}{\mathbb{P}} 
\newcommand{\E}{\mathbb{E}}

\newcommand{\given}{\,\middle|\,}
\newcommand{\ind}{\mathbf{1}}

% Common “defined as” symbol
\newcommand{\defeq}{\vcentcolon=}

% =========================
% 8) Problem command (TOC + PDF bookmarks)
% Usage:
%   \problem{<label>}
%   \problem[Short title]{<label>}
% Example:
%   \problem{1.4}
%   \problem[Sample space]{1.4}
% =========================
\makeatletter
\newcommand{\problem}[2][]{%
  \def\prob@title{Problem #2\if\relax\detokenize{#1}\relax\else\ (\textit{#1})\fi}%
  \phantomsection%
  \pdfbookmark[1]{\prob@title}{prob#2}%
  \addcontentsline{toc}{section}{\prob@title}%
  \section*{\prob@title}%
  \vspace{-0.25em}\hrule\vspace{0.8em}%
}
\makeatother

% =========================
% 9) Solution environment
% =========================
\newenvironment{solution}{%
  \noindent\textbf{Solution.}\ \ignorespaces
}{\par}

% Optional scratch / planning block
\newenvironment{scratch}{%
  \begin{boxnote}\textbf{Scratch / Setup.}\ \small
}{\end{boxnote}}

% Final answer command: expects math only (no $$ or \[ \])
\newcommand{\finalanswer}[1]{%
  \par\medskip
  \noindent\textbf{Final:}\ \fbox{$#1$}\par
} % AGAIN: MATH ONLY INSIDE THE BOX; if word answer, use \text{} inside or use \fbox{TEXT} or \framebox

% =========================
% 10) Title block
% =========================
\title{Homework \HomeworkNumber}
\author{\YourName \\
\CourseCode\ --- \CourseTitle \\
\CourseSection \\
Instructor: \Instructor}
\date{Due: \DueDate}

\begin{document}
\maketitle

% \section*{Collaboration / Resources}
% (Optional) Write “I worked alone” or list collaborators/resources.

\tableofcontents
\bigskip
\hrule
\bigskip
\newpage

% ============================================================
% Problems (duplicate blocks as needed)
% ============================================================


\problem{4.34}
A taxi company has a large fleet of cars. On average, there are 3 accidents each week. What is the probability that at most 2 accidents happen next week? Make some reasonable assumption in order to be able to answer the question. Simplify your answer as much as possible.

\begin{solution}
\begin{scratch}
Let $X$ be the number of accidents next week. A standard modeling assumption is:
each taxi has a small probability of an accident during the week, accidents for
different taxis are approximately independent, and the fleet is large. Then a
Poisson model is appropriate, with parameter equal to the weekly mean:
\[
X\sim \mathrm{Poisson}(3).
\]
\end{scratch}

Assuming accidents occur as independent rare events across a large fleet, the
number of accidents in one week is well modeled by a Poisson random variable
with mean $3$. Thus
\[
P(X\le 2)=P(X=0)+P(X=1)+P(X=2).
\]
Using the Poisson p.m.f.,
\[
P(X\le 2)=e^{-3}\left(\frac{3^0}{0!}+\frac{3^1}{1!}+\frac{3^2}{2!}\right)
=e^{-3}\left(1+3+\frac{9}{2}\right)
=\frac{17}{2}e^{-3}.
\]

\finalanswer{P(X\le 2)=\frac{17}{2}e^{-3}}

\end{solution}

\newpage

\problem{4.42}
On average 20\% of the gadgets produced by a factory are mildly defective. I buy a box of 100 gadgets. Assume this is a random sample from the production of the factory. Let $A$ be the event that less than 15 gadgets in the random sample of 100 are mildly defective.
\begin{enumerate}[label=(\alph*)]
\item Give an exact expression for $P(A)$, without attempting to evaluate it.
\item Use either the normal or the Poisson approximation, whichever is appropriate, to give an approximation of $P(A)$.
\end{enumerate}

\begin{solution}
\begin{scratch}
Let $X$ be the number of mildly defective gadgets in the box. Then
\[
X\sim \mathrm{Bin}(100,0.2),
\qquad
A=\{X<15\}=\{X\le 14\}.
\]
For approximation, compare Poisson and normal:
\[
np=20,\qquad np(1-p)=16.
\]
The normal approximation is the better choice here. Using continuity correction,
\[
P(X\le 14)\approx P\!\left(Z\le \frac{14.5-20}{4}\right)=P(Z\le -1.375).
\]
\end{scratch}

\begin{enumerate}[label=(\alph*)]
\item Since $X\sim\mathrm{Bin}(100,0.2)$,
\[
P(A)=P(X\le 14)=\sum_{k=0}^{14}\binom{100}{k}(0.2)^k(0.8)^{100-k}.
\]
\item Because the mean $np=20$ and variance $np(1-p)=16$ are both reasonably
large, the normal approximation is more appropriate than the Poisson
approximation. Approximating $X$ by $N(20,16)$ and using continuity
correction,
\[
P(A)=P(X\le 14)\approx P\!\left(Z\le \frac{14.5-20}{4}\right)
=P(Z\le -1.375).
\]
Numerically,
\[
P(A)\approx \Phi(-1.375)\approx 0.0846.
\]
\end{enumerate}

\finalanswer{\text{(a)}\ \sum_{k=0}^{14}\binom{100}{k}(0.2)^k(0.8)^{100-k},\qquad
\text{(b)}\ P(A)\approx \Phi(-1.375)\approx 0.0846}

\end{solution}

\newpage

\problem{4.46}
Jessica flips a fair coin 5 times every morning, for 30 days straight. Let $X$ be the number of mornings over these 30 days on which all 5 flips are tails. Use either the normal or the Poisson approximation, whichever is more appropriate, to give an estimate for the probability $P(X=2)$. Justify your choice of approximation.

\begin{solution}
\begin{scratch}
For one morning, the probability of getting 5 tails in 5 flips is
\[
p=\left(\frac12\right)^5=\frac{1}{32}.
\]
If $X$ is the number of such mornings in 30 days, then
\[
X\sim \mathrm{Bin}\left(30,\frac{1}{32}\right).
\]
Here
\[
np=\frac{15}{16},\qquad np^2=\frac{15}{512},
\]
so the event is rare and the Poisson approximation is more appropriate than the
normal approximation.
\end{scratch}

Let $X$ be the number of mornings, out of 30, on which all 5 flips are tails.
Then
\[
X\sim \mathrm{Bin}\left(30,\frac{1}{32}\right).
\]
Since $p=1/32$ is small and $np^2=15/512$ is very small, the Poisson
approximation is appropriate. We approximate $X$ by
\[
Y\sim \mathrm{Poisson}\left(\lambda=\frac{30}{32}\right)
=\mathrm{Poisson}\left(\frac{15}{16}\right).
\]
Therefore,
\[
P(X=2)\approx P(Y=2)
=e^{-15/16}\frac{(15/16)^2}{2!}
=\frac{225}{512}e^{-15/16}.
\]
Numerically,
\[
P(X=2)\approx 0.1721.
\]

\finalanswer{P(X=2)\approx \frac{225}{512}e^{-15/16}\approx 0.1721}

\end{solution}

\newpage

\problem{5.2}
Suppose that $X$ has moment generating function
\[
M_X(t)=\frac{1}{2}+\frac{1}{3}e^{-4t}+\frac{1}{6}e^{5t}.
\]
\begin{enumerate}[label=(\alph*)]
\item Find the mean and variance of $X$ by differentiating the moment generating function to find moments.
\item Find the probability mass function of $X$. Use the probability mass function to check your answer for part (a).
\end{enumerate}

\begin{solution}
\begin{scratch}
Differentiate the mgf to get moments:
\[
M_X'(t)=-\frac{4}{3}e^{-4t}+\frac{5}{6}e^{5t},
\qquad
M_X''(t)=\frac{16}{3}e^{-4t}+\frac{25}{6}e^{5t}.
\]
Also, because
\[
M_X(t)=\sum_x P(X=x)e^{tx},
\]
the exponential terms show the support and probabilities directly.
\end{scratch}

\begin{enumerate}[label=(\alph*)]
\item Using the mgf,
\[
E[X]=M_X'(0)=-\frac{4}{3}+\frac{5}{6}=-\frac{8}{6}+\frac{5}{6}=-\frac{1}{2}.
\]
Also,
\[
E[X^2]=M_X''(0)=\frac{16}{3}+\frac{25}{6}
=\frac{32}{6}+\frac{25}{6}=\frac{57}{6}=\frac{19}{2}.
\]
Hence
\[
\operatorname{Var}(X)=E[X^2]-\bigl(E[X]\bigr)^2
=\frac{19}{2}-\left(-\frac{1}{2}\right)^2
=\frac{19}{2}-\frac{1}{4}
=\frac{37}{4}.
\]
\item Since
\[
M_X(t)=\frac12+\frac13e^{-4t}+\frac16e^{5t},
\]
we read off the distribution:
\[
P(X=0)=\frac12,\qquad P(X=-4)=\frac13,\qquad P(X=5)=\frac16.
\]
Checking part (a) from the p.m.f.,
\[
E[X]=(-4)\cdot \frac13+0\cdot \frac12+5\cdot \frac16=-\frac{4}{3}+\frac{5}{6}=-\frac12,
\]
which agrees with the mgf calculation.
\end{enumerate}

\finalanswer{E[X]=-\frac12,\quad \operatorname{Var}(X)=\frac{37}{4},\quad
P(X=-4)=\frac13,\ P(X=0)=\frac12,\ P(X=5)=\frac16}

\end{solution}

\newpage

\problem{5.8}
Let $X\sim\mathrm{Unif}[-1,2]$. Find the probability density function of the random variable $Y=X^2$.

\begin{solution}
\begin{scratch}
Let $Y=X^2$. Since $X\sim\mathrm{Unif}[-1,2]$, its density is
\[
f_X(x)=\frac13\quad \text{for }-1\le x\le 2.
\]
The map $x\mapsto x^2$ is not one-to-one:
for $0<y<1$, the preimages are $\pm\sqrt y$;
for $1<y<4$, only $\sqrt y$ lies in $[-1,2]$.
Use
\[
f_Y(y)=\sum_{x:\,x^2=y} f_X(x)\frac{1}{|2x|}.
\]
\end{scratch}

The possible values of $Y=X^2$ are $0\le Y\le 4$.

If $0<y<1$, both $\sqrt y$ and $-\sqrt y$ belong to the support of $X$, so
\[
f_Y(y)=f_X(\sqrt y)\frac{1}{2\sqrt y}+f_X(-\sqrt y)\frac{1}{2\sqrt y}
=\frac13\cdot\frac{1}{2\sqrt y}+\frac13\cdot\frac{1}{2\sqrt y}
=\frac{1}{3\sqrt y}.
\]
If $1<y<4$, only $\sqrt y$ belongs to $[-1,2]$, so
\[
f_Y(y)=f_X(\sqrt y)\frac{1}{2\sqrt y}
=\frac13\cdot\frac{1}{2\sqrt y}
=\frac{1}{6\sqrt y}.
\]
Outside $(0,4)$ the density is $0$. Therefore,
\[
f_Y(y)=
\begin{cases}
\dfrac{1}{3\sqrt y}, & 0<y<1,\\[6pt]
\dfrac{1}{6\sqrt y}, & 1<y<4,\\[6pt]
0, & \text{otherwise.}
\end{cases}
\]
The value at $y=1$ is irrelevant for the density.

\finalanswer{f_Y(y)=
\begin{cases}
\dfrac{1}{3\sqrt y}, & 0<y<1,\\[4pt]
\dfrac{1}{6\sqrt y}, & 1<y<4,\\[4pt]
0, & \text{otherwise}
\end{cases}}

\end{solution}

\newpage

\problem{5.10}
Suppose that $X$ has moment generating function
\[
M(t)=\left(\frac{1}{5}+\frac{4}{5}e^t\right)^{30}.
\]
What is the distribution of $X$?

\emph{Hint.} Do Exercise 5.9 first.

\begin{solution}
\begin{scratch}
From Exercise 5.9, the mgf of a binomial random variable is
\[
M(t)=(1-p+pe^t)^n.
\]
Compare this with
\[
\left(\frac15+\frac45 e^t\right)^{30}.
\]
\end{scratch}

By the binomial mgf formula,
\[
M_X(t)=(1-p+pe^t)^n
\]
for $X\sim\mathrm{Bin}(n,p)$. Since
\[
M(t)=\left(\frac15+\frac45 e^t\right)^{30},
\]
we identify
\[
n=30,\qquad p=\frac45.
\]
Therefore,
\[
X\sim \mathrm{Bin}\left(30,\frac45\right).
\]
Equivalently, its p.m.f. is
\[
P(X=k)=\binom{30}{k}\left(\frac45\right)^k\left(\frac15\right)^{30-k},
\qquad k=0,1,\dots,30.
\]

\finalanswer{X\sim \mathrm{Bin}\left(30,\frac45\right)}

\end{solution}

\newpage

\problem{5.18}
Let $X\sim\mathrm{Geom}(p)$.
\begin{enumerate}[label=(\alph*)]
\item Compute the moment generating function $M_X(t)$ of $X$. Be careful about the possibility that $M_X(t)$ might be infinite.
\item Use the moment generating function to compute the mean and the variance of $X$.
\end{enumerate}

\begin{solution}
\begin{scratch}
Write $q=1-p$. Then for $X\sim\mathrm{Geom}(p)$ on $\{1,2,\dots\}$,
\[
P(X=k)=pq^{k-1}.
\]
The mgf becomes a geometric series:
\[
M_X(t)=\sum_{k=1}^\infty e^{tk}pq^{k-1}
=pe^t\sum_{k=0}^\infty (qe^t)^k.
\]
\end{scratch}

\begin{enumerate}[label=(\alph*)]
\item For those $t$ such that $qe^t<1$, the series converges and
\[
M_X(t)=pe^t\sum_{k=0}^\infty (qe^t)^k
=\frac{pe^t}{1-qe^t}
=\frac{pe^t}{1-(1-p)e^t}.
\]
Thus
\[
M_X(t)=
\begin{cases}
\dfrac{pe^t}{1-(1-p)e^t}, & t<\ln\!\left(\dfrac{1}{1-p}\right),\\[8pt]
\infty, & t\ge \ln\!\left(\dfrac{1}{1-p}\right).
\end{cases}
\]
\item Differentiate
\[
M_X(t)=\frac{pe^t}{1-qe^t}.
\]
Then
\[
M_X'(t)=\frac{pe^t}{(1-qe^t)^2},
\]
so
\[
E[X]=M_X'(0)=\frac{p}{(1-q)^2}=\frac{p}{p^2}=\frac{1}{p}.
\]
Differentiating once more,
\[
M_X''(t)=\frac{pe^t}{(1-qe^t)^2}+\frac{2pq\,e^{2t}}{(1-qe^t)^3},
\]
hence
\[
E[X^2]=M_X''(0)=\frac{1}{p}+\frac{2(1-p)}{p^2}
=\frac{2-p}{p^2}.
\]
Therefore,
\[
\operatorname{Var}(X)=E[X^2]-\bigl(E[X]\bigr)^2
=\frac{2-p}{p^2}-\frac{1}{p^2}
=\frac{1-p}{p^2}.
\]
\end{enumerate}

\finalanswer{M_X(t)=\frac{pe^t}{1-(1-p)e^t}\ \text{for }t<\ln\!\left(\frac{1}{1-p}\right),\quad
E[X]=\frac{1}{p},\quad \operatorname{Var}(X)=\frac{1-p}{p^2}}

\end{solution}

\newpage

\problem{5.20}
Suppose that random variable $X$ has density function $f(x)=\frac{1}{2}e^{-|x|}$.
\begin{enumerate}[label=(\alph*)]
\item Compute the moment generating function $M_X(t)$ of $X$. Be careful about the possibility that $M_X(t)$ might be infinite.
\item Use the moment generating function to compute the $n$th moment of $X$.
\end{enumerate}

\begin{solution}
\begin{scratch}
Split the integral at $0$. For $x\ge 0$, $|x|=x$; for $x<0$, $|x|=-x$.
So
\[
M_X(t)=\frac12\int_0^\infty e^{tx}e^{-x}\,dx
\;+\;\frac12\int_{-\infty}^0 e^{tx}e^{x}\,dx.
\]
Both integrals converge exactly when $-1<t<1$.
\end{scratch}

\begin{enumerate}[label=(\alph*)]
\item We compute
\[
M_X(t)=\frac12\int_0^\infty e^{-(1-t)x}\,dx+\frac12\int_{-\infty}^0 e^{(1+t)x}\,dx.
\]
The first integral converges when $t<1$ and equals $\dfrac{1}{2(1-t)}$.
The second converges when $t>-1$ and equals $\dfrac{1}{2(1+t)}$.
Hence
\[
M_X(t)=
\begin{cases}
\dfrac{1}{2(1-t)}+\dfrac{1}{2(1+t)}=\dfrac{1}{1-t^2}, & -1<t<1,\\[8pt]
\infty, & |t|\ge 1.
\end{cases}
\]
\item For $|t|<1$,
\[
M_X(t)=\frac{1}{1-t^2}=\sum_{k=0}^\infty t^{2k}.
\]
Comparing with
\[
M_X(t)=\sum_{n=0}^\infty \frac{E[X^n]}{n!}t^n,
\]
we see that odd powers do not appear, so
\[
E[X^n]=0 \qquad \text{for odd }n.
\]
If $n=2k$ is even, the coefficient of $t^{2k}$ is $1$, so
\[
\frac{E[X^{2k}]}{(2k)!}=1,
\qquad\text{hence}\qquad
E[X^{2k}]=(2k)!.
\]
\end{enumerate}

\finalanswer{M_X(t)=\frac{1}{1-t^2}\ \text{for }|t|<1,\qquad
E[X^n]=
\begin{cases}
0, & n\ \text{odd},\\
n!, & n\ \text{even}
\end{cases}}

\end{solution}

\newpage

\problem{5.24}
Let $X\sim N(0,1)$ and $Y=e^X$. $Y$ is called a \emph{log-normal} random variable.
\begin{enumerate}[label=(\alph*)]
\item Find the probability density function of $Y$.
\item Find the $n$th moment $E(Y^n)$ of $Y$.
\end{enumerate}

\emph{Hint.} Do not compute the moment generating function of $Y$. Instead relate the $n$th moment of $Y$ to an expectation of $X$ that you know.

\begin{solution}
\begin{scratch}
Since $Y=e^X$ and the exponential function is increasing, use the cdf method:
\[
F_Y(y)=P(Y\le y)=P(e^X\le y).
\]
Also $Y>0$ always. For moments,
\[
Y^n=e^{nX},
\]
so $E[Y^n]$ is the standard normal mgf evaluated at $n$.
\end{scratch}

\begin{enumerate}[label=(\alph*)]
\item Because $Y=e^X>0$, we have $F_Y(y)=0$ for $y\le 0$. For $y>0$,
\[
F_Y(y)=P(e^X\le y)=P(X\le \ln y)=\Phi(\ln y),
\]
where $\Phi$ is the standard normal cdf. Differentiating,
\[
f_Y(y)=\Phi'(\ln y)\cdot \frac{1}{y}
=\frac{1}{\sqrt{2\pi}}e^{-(\ln y)^2/2}\cdot \frac{1}{y},
\qquad y>0.
\]
Thus
\[
f_Y(y)=
\begin{cases}
\dfrac{1}{y\sqrt{2\pi}}e^{-(\ln y)^2/2}, & y>0,\\[8pt]
0, & y\le 0.
\end{cases}
\]
\item Since $Y=e^X$,
\[
E[Y^n]=E[e^{nX}]=M_X(n).
\]
For $X\sim N(0,1)$,
\[
M_X(t)=e^{t^2/2},
\]
so
\[
E[Y^n]=e^{n^2/2}.
\]
\end{enumerate}

\finalanswer{f_Y(y)=
\begin{cases}
\dfrac{1}{y\sqrt{2\pi}}e^{-(\ln y)^2/2}, & y>0,\\[6pt]
0, & y\le 0
\end{cases},
\qquad E[Y^n]=e^{n^2/2}}

\end{solution}

\newpage

\problem{5.32}
Suppose that $X$ is a uniform random variable on the interval $(0,1)$ and let $Y=1/X$. Find the probability density function of $Y$.

\begin{solution}
\begin{scratch}
Use the cdf. Since $0<X<1$, the transformed variable is
\[
Y=\frac{1}{X}>1.
\]
For $y>1$,
\[
F_Y(y)=P\!\left(\frac1X\le y\right)=P\!\left(X\ge \frac1y\right).
\]
\end{scratch}

Because $X$ is uniform on $(0,1)$, its density is $1$ on $(0,1)$. Since
$Y=1/X$, the support of $Y$ is $(1,\infty)$.

If $y\le 1$, then $F_Y(y)=0$. If $y>1$, then
\[
F_Y(y)=P\!\left(\frac1X\le y\right)
=P\!\left(X\ge \frac1y\right)
=1-\frac1y.
\]
Differentiating for $y>1$ gives
\[
f_Y(y)=\frac{d}{dy}\left(1-\frac1y\right)=\frac{1}{y^2}.
\]
Therefore,
\[
f_Y(y)=
\begin{cases}
\dfrac{1}{y^2}, & y>1,\\[6pt]
0, & y\le 1.
\end{cases}
\]

\finalanswer{f_Y(y)=
\begin{cases}
\dfrac{1}{y^2}, & y>1,\\[4pt]
0, & y\le 1
\end{cases}}

\end{solution}


\end{document}
